\documentclass[11pt]{article}

% ============================================================
% GA-PRAC-01 — Comparación entre modelos de proceso de requisitos en escenarios reales
% Ingeniería de Requerimientos (ISR-401) · UTEQ · SPA 2026-2027
% Compilación: pdflatex -> bibtex -> pdflatex -> pdflatex
% ============================================================

\usepackage[utf8]{inputenc}
\usepackage[T1]{fontenc}
% babel no disponible en este entorno
\usepackage[scaled]{helvet}
\renewcommand{\familydefault}{\sfdefault}
\usepackage[a4paper,top=2.6cm,bottom=2.3cm,left=2cm,right=2cm,headheight=34pt]{geometry}
\usepackage{graphicx}
\usepackage{array}
\usepackage{tabularx}
\usepackage{longtable}
\usepackage{pdflscape}
\usepackage{colortbl}
\usepackage{xcolor}
\usepackage{enumitem}
\usepackage{fancyhdr}
\usepackage{lastpage}
\usepackage[most]{tcolorbox}
\usepackage{amsmath}
\usepackage{amssymb}
\usepackage[numbers,sort&compress]{natbib}
\usepackage{url}
\usepackage{needspace}
\usepackage{caption}
\usepackage[hidelinks]{hyperref}

% ---------- Textos en español (sin babel) ----------
\renewcommand{\tablename}{Tabla}
\renewcommand{\refname}{Referencias}

% ---------- Paleta institucional ----------
\definecolor{utegreen}{HTML}{0B7A3B}
\definecolor{utegreendk}{HTML}{075C2C}
\definecolor{midblue}{HTML}{1F5C8B}
\definecolor{gold}{HTML}{D4A017}
\definecolor{rowgray}{HTML}{EFEFEF}
\definecolor{boxgray}{HTML}{F4F6F4}
\definecolor{lightred}{HTML}{FBEAEA}
\definecolor{lightgold}{HTML}{FBF3DC}
\definecolor{lightblue}{HTML}{EAF1F7}
\definecolor{lightgreen}{HTML}{DCEFE3}

% ---------- Columnas personalizadas ----------
\newcolumntype{L}[1]{>{\raggedright\arraybackslash}p{#1}}
\newcolumntype{C}[1]{>{\centering\arraybackslash}p{#1}}
\newcolumntype{Y}{>{\raggedright\arraybackslash}X}

% ---------- Encabezado y pie ----------
\pagestyle{fancy}
\fancyhf{}
\lhead{\footnotesize UTEQ \,$\cdot$\, Ingeniería de Requerimientos (ISR-401)}
\rhead{\footnotesize GA-PRAC-01 \,$\cdot$\, SPA 2026--2027}
\lfoot{\footnotesize Prof. Gleiston C. Guerrero-Ulloa, PhD}
\rfoot{\footnotesize Página \thepage\ de \pageref{LastPage}}
\renewcommand{\headrulewidth}{0.6pt}
\renewcommand{\footrulewidth}{0.4pt}

\setlength{\parindent}{0pt}
\setlength{\parskip}{5pt}
\setlength{\emergencystretch}{3em}
\renewcommand{\arraystretch}{1.25}

% ---------- Atajos ----------
\newcommand{\thdr}[1]{\cellcolor{utegreen}\textcolor{white}{\textbf{#1}}}
\newcommand{\seccion}[1]{\section{\textcolor{utegreendk}{#1}}}
\newcommand{\subseccion}[1]{\subsection{\textcolor{midblue}{#1}}}
\renewcommand{\thesection}{\arabic{section}.}
\renewcommand{\thesubsection}{\arabic{section}.\arabic{subsection}.}

\newtcolorbox{cajaclave}[1][]{colback=boxgray,colframe=utegreen,boxrule=0.8pt,arc=2mm,fonttitle=\bfseries,breakable,#1}
\newtcolorbox{cajapiso}[1][]{colback=lightred,colframe=red!60!black,boxrule=1pt,arc=2mm,fonttitle=\bfseries,breakable,#1}
\newtcolorbox{cajaaviso}[1][]{colback=lightgold,colframe=gold!80!black,boxrule=0.8pt,arc=2mm,fonttitle=\bfseries,breakable,#1}

\begin{document}

\noindent\includegraphics[width=\linewidth]{logoband.png}

\vspace{4pt}
\begin{center}
	{\small Facultad de Ciencias de la Computación \,\textcolor{gold}{$\bullet$}\, Carrera de Software \,\textcolor{gold}{$\bullet$}\, Ingeniería de Requerimientos (ISR-401)}\\[8pt]
	{\LARGE\textbf{\textcolor{utegreendk}{GUÍA Y RÚBRICA DE EVALUACIÓN}}}\\[4pt]
	{\large\textbf{\textcolor{midblue}{ACTIVIDAD FORMATIVA}}}\\[6pt]
	{\Large\textbf{GA-PRAC-01}}\\[2pt]
	{\large Comparación entre modelos de proceso de requisitos en escenarios reales}
\end{center}

\vspace{2pt}
\noindent\textcolor{gold}{\rule{\linewidth}{1.5pt}}

\seccion{Datos generales}

\noindent\begin{tabularx}{\linewidth}{|L{0.30\linewidth}|Y|}
	\hline
	\thdr{Campo} & \thdr{Detalle}\\
	\hline
	Asignatura & Ingeniería de Requerimientos (ISR-401) \,$\cdot$\, 4.\textsuperscript{o} nivel \,$\cdot$\, Paralelo A\\
	\hline
	\rowcolor{rowgray} Docente & Gleiston Cicerón Guerrero Ulloa, PhD \,$\cdot$\, \texttt{gguerrero@uteq.edu.ec}\\
	\hline
	Unidad & Unidad 1. Fundamentos de Ingeniería de Requisitos y Proceso de Requisitos\\
	\hline
	\rowcolor{rowgray} Resultado de aprendizaje & Aplica los fundamentos teóricos de la IR identificando tipos, propiedades y el proceso de requisitos para analizar su aplicabilidad en proyectos de software reales.\\
	\hline
	Carácter & Formativo: sin calificación; su entrega válida es requisito de la sumativa vinculada\\
	\hline
	\rowcolor{rowgray} Prepara (sumativas vinculadas) & GA-PE-U, EV-PRAC-01\\
	\hline
	Modalidad & Individual\\
	\hline
	\rowcolor{rowgray} Semana & 2 (2 de noviembre -- 8 de noviembre de 2026)\\
	\hline
	Publicación de la guía & lunes 2 de noviembre de 2026\\
	\hline
	\rowcolor{rowgray} Cierre & sábado 7 de noviembre de 2026, a las 23:59\\
	\hline
	Repositorio & Repositorio personal \texttt{ISR401-<ApellidoNombre>}\\
	\hline
\end{tabularx}

\seccion{Propósito}

Comparar los modelos de proceso de requisitos en escenarios reales para aprender a elegir el más adecuado según el tipo de sistema, la estabilidad de los requisitos y la disponibilidad de las partes interesadas.

\seccion{Fundamento}

El proceso de requisitos puede organizarse de forma secuencial, iterativa o ágil, y su elección depende de factores como la volatilidad de los requisitos, el tipo de sistema y la relación con el cliente \citep{pohl_2025,swebok_v4_2024,ireb_cpre_fl_2026}.

\seccion{Desarrollo de la actividad}

\begin{enumerate}[leftmargin=20pt,itemsep=4pt]
	\item \textbf{Escenarios.} Lea los 4 escenarios del Anexo A (sistema embebido de una planta, aplicación móvil de un emprendimiento, sistema de un GAD con contrato público y plataforma de una cooperativa).
	\item \textbf{Criterios.} Defina al menos 5 criterios de comparación (estabilidad, criticidad, disponibilidad del cliente, tamaño, regulación).
	\item \textbf{Matriz.} Elabore la matriz escenario $\times$ modelo con la valoración de cada criterio y elija un modelo por escenario.
	\item \textbf{Framework.} Para uno de los escenarios, describa qué actividades centrales y transversales del framework de IR se ejecutan en cada iteración o fase.
\end{enumerate}

\seccion{Productos entregables}

{\small
\begin{longtable}{|L{0.27\linewidth}|L{0.43\linewidth}|L{0.21\linewidth}|}
	\caption{Productos que deben estar en el repositorio al cierre.}\\
	\hline
	\thdr{Ruta en el repositorio} & \thdr{Contenido} & \thdr{Mínimo exigido}\\
	\hline
	\endfirsthead
	\hline
	\thdr{Ruta en el repositorio} & \thdr{Contenido} & \thdr{Mínimo exigido}\\
	\hline
	\endhead
	\texttt{GA-PRAC-01/\allowbreak{}main.tex y main.pdf} & Criterios, matriz, elección justificada y mapeo al framework & 4 escenarios\\
	\hline
\end{longtable}
}

\seccion{Normas de entrega}

\begin{enumerate}[leftmargin=20pt,itemsep=3pt]
	\item \textbf{Repositorio.} Todo el trabajo se desarrolla en el repositorio personal \texttt{ISR401-<ApellidoNombre>}, que debe ser público en GitHub. Allí deben estar los fuentes (\texttt{.tex}, \texttt{.bib}, figuras, fuentes editables de los modelos y evidencias) y el PDF compilado.
	\item \textbf{Documento en el SGA.} En la tarea del SGA se sube solamente un PDF de \textbf{una sola página} con los datos informativos: estudiante o estudiantes, código y nombre de la tarea, fecha límite de entrega y la URL del repositorio escrita completa en una sola línea. No se sube ningún otro archivo.
	\item \textbf{\LaTeX{} reproducible.} El PDF debe poder regenerarse clonando el repositorio y compilando el archivo \texttt{.tex}. El \texttt{README.md} documenta compilador, archivo principal, orden de comandos y dependencias.
	\item \textbf{Referencias.} Estilo IEEE. Toda referencia debe existir y sus datos (autores, año, título, revista o editorial, volumen, número, páginas y DOI o ISBN) deben corresponder exactamente a la obra real.
	\item \textbf{Plazo.} El cierre es el sábado 7 de noviembre de 2026, a las 23:59. Se evalúa el contenido del repositorio hasta ese momento; lo agregado después no se considera.
	\item \textbf{Trabajo anticipado.} El trabajo se realiza durante la semana, no a última hora; el historial de commits debe mostrarlo. La lista de cotejo registra la constancia observada.
	\item \textbf{Datos personales.} Ningún dato personal de terceros (cédulas, teléfonos, correos personales, audios o videos de participantes, consentimientos firmados) se publica en el repositorio. Esas evidencias se guardan en una carpeta institucional privada (cuenta \texttt{@uteq.edu.ec}) compartida con el docente, y en el repositorio se registra un manifiesto con nombre de archivo, fecha, duración y huella SHA-256 \citep{lopdp_2021}.
	\item \textbf{Uso de IA como apoyo.} Si se usa una herramienta de IA solo como apoyo (ortografía, error de compilación), se declara en una sección del documento: herramienta, versión, propósito y cómo se verificó lo obtenido.
\end{enumerate}

\seccion{Condiciones de entrega válida}

\begin{cajapiso}[title={Condiciones de entrega válida: su incumplimiento equivale a no entregar}]
	\begin{enumerate}[label=\textbf{P\arabic*.},leftmargin=24pt,itemsep=3pt]
		\item \textbf{PDF de una sola página con datos informativos y URL.} En el SGA se sube solamente un PDF de UNA SOLA página con los datos informativos (estudiante o estudiantes, tarea, fecha límite de entrega) y, en una sola línea, la URL del repositorio. Si el PDF tiene más de una página, si se suben otros archivos, si la URL no cabe en una sola línea o está rota, o si el repositorio no es accesible, la actividad formativa se considera \textbf{NO ENTREGADA} y el estudiante no puede presentar la sumativa vinculada, en la que obtiene CERO.
		\item \textbf{PDF reproducible.} Si el PDF no se genera correctamente a partir del \LaTeX{} al clonar el repositorio y compilar el archivo \texttt{.tex} siguiendo el \texttt{README.md} (compilador, orden de comandos, archivo principal y dependencias), la actividad formativa se considera \textbf{NO ENTREGADA} y el estudiante no puede presentar la sumativa vinculada, en la que obtiene CERO.
		\item \textbf{Plazo.} Si el contenido evaluable se sube después de la fecha y hora de cierre, la actividad formativa se considera \textbf{NO ENTREGADA} y el estudiante no puede presentar la sumativa vinculada, en la que obtiene CERO.
		\item \textbf{Integridad.} Plagio verificado, cualquiera de los indicios graves G1 a G6 de uso de IA, o contenido generado por IA presentado como propio: la actividad formativa se considera \textbf{NO ENTREGADA} y el estudiante no puede presentar la sumativa vinculada, en la que obtiene CERO, sin perjuicio de las acciones disciplinarias que correspondan.
	\end{enumerate}
\end{cajapiso}

\seccion{Integridad académica e indicios de uso de IA}

Todo trabajo debe ser de autoría propia. Se aplica el Código de Ética de la ACM \citep{acm_code_2018} y la sección 6.1 del encuadre.

\begin{cajapiso}[title={Indicios graves de uso de IA: cada uno, por sí solo, implica CERO (P4)}]
	\begin{enumerate}[label=\textbf{G\arabic*.},leftmargin=24pt,itemsep=3pt]
		\item \textbf{Referencias con datos que no corresponden a la obra real:} DOI que lleva a una obra con autores, título, año, revista, volumen o páginas diferentes; DOI, ISBN o URL inexistentes; obras que no existen.
		\item \textbf{Citas que la fuente no respalda:} afirmaciones, cifras o citas textuales atribuidas a una fuente que no las contiene, o cláusulas, secciones o tablas de una norma (IEEE 29148, ISO/\allowbreak{}IEC 25010) que no existen o no dicen lo afirmado.
		\item \textbf{Restos de un asistente de IA} en el documento, el código, los comentarios o los commits (por ejemplo, «Claro, aquí tienes\ldots», «Como modelo de lenguaje\ldots»).
		\item \textbf{Evidencias de campo fabricadas:} entrevistas, encuestas, participantes, transcripciones, actas o fotografías que no existen o no tienen respaldo original.
		\item \textbf{Resultados o artefactos inexistentes:} el texto describe modelos, matrices, pruebas o métricas que no están en el repositorio o no coinciden con él.
		\item \textbf{Incapacidad de explicar lo entregado:} el estudiante no puede justificar ni modificar, cuando se le pide, lo que escribió, modeló o especificó.
	\end{enumerate}
\end{cajapiso}

\begin{cajaaviso}[title={Indicios complementarios (se acumulan)}]
	\begin{enumerate}[label=\textbf{I\arabic*.},leftmargin=24pt,itemsep=3pt]
		\item Redacción uniforme de principio a fin, o estilo distinto al del estudiante en otras entregas o en su exposición oral.
		\item Uso excesivo de los mismos conectores y fórmulas («además», «cabe destacar», «es importante señalar», «en conclusión»).
		\item Contenido genérico: requisitos, actores o reglas aplicables a cualquier sistema, sin relación con lo elicitado.
		\item Entrega en uno o pocos commits masivos, o mensajes de commit idénticos o automáticos.
		\item Residuos de formato ajeno a \LaTeX{} (marcas de Markdown, emojis) o metadatos del PDF de otra herramienta.
		\item Estructura, redacción o errores muy similares entre estudiantes o equipos distintos.
		\item Traducción literal o mezcla de idiomas sin justificación.
	\end{enumerate}
	Con dos indicios complementarios se resta 1,5 puntos (nivel B) y con tres o más se restan 3,0 (nivel C). La declaración de uso de IA no exime de responsabilidad ni anula estos indicios.
\end{cajaaviso}

\seccion{Lista de cotejo para la retroalimentación}

La actividad no tiene calificación. El docente marca cada indicador y escribe una observación que el estudiante debe atender en la actividad sumativa vinculada. \textbf{L} = Logrado (cumple con la calidad descrita); \textbf{EP} = En proceso (presente pero incompleto, impreciso o con errores que se señalan); \textbf{NL} = No logrado (ausente, incorrecto o no verificable). Que un elemento esté presente no basta: se valora su calidad.

{\small
\begin{longtable}{|C{0.04\linewidth}|L{0.47\linewidth}|C{0.035\linewidth}|C{0.035\linewidth}|C{0.035\linewidth}|L{0.17\linewidth}|}
	\caption{Lista de cotejo de GA-PRAC-01.}\\
	\hline
	\thdr{N.\textsuperscript{o}} & \thdr{Indicador y calidad esperada (Logrado)} & \thdr{L} & \thdr{EP} & \thdr{NL} & \thdr{Observación}\\
	\hline
	\endfirsthead
	\hline
	\thdr{N.\textsuperscript{o}} & \thdr{Indicador y calidad esperada (Logrado)} & \thdr{L} & \thdr{EP} & \thdr{NL} & \thdr{Observación}\\
	\hline
	\endhead
	1 & Al menos 5 criterios definidos con significado claro. & $\square$ & $\square$ & $\square$ & \\
	\hline
	2 & Matriz completa para los 4 escenarios y 3 modelos. & $\square$ & $\square$ & $\square$ & \\
	\hline
	3 & Elección justificada con los criterios, no por preferencia. & $\square$ & $\square$ & $\square$ & \\
	\hline
	4 & Actividades centrales y transversales del framework correctamente ubicadas. & $\square$ & $\square$ & $\square$ & \\
	\hline
	5 & Documento propio, compilado y en el repositorio personal. & $\square$ & $\square$ & $\square$ & \\
	\hline
\end{longtable}
}

\begin{cajaclave}[title={Cómo se usa esta retroalimentación}]
	El docente devuelve la lista de cotejo en el SGA antes de la sesión siguiente. Cada indicador marcado como EP o NL debe corregirse antes de presentar GA-PE-U y EV-PRAC-01; el docente verificará esa corrección al calificarla. Entregar esta formativa de forma válida es requisito para presentar GA-PE-U y EV-PRAC-01.
\end{cajaclave}

% ============================================================
% ANEXOS: material que entrega el docente
% ============================================================
\clearpage
\seccion{Anexo A. Escenarios para la comparación}

Los cuatro escenarios de este anexo describen organizaciones, personas y datos ficticios creados para esta actividad; cualquier parecido con casos reales es coincidencia. Lea cada escenario completo antes de llenar la matriz y subraye los datos que le sirvan como evidencia para cada criterio de comparación. Los escenarios no tienen una única respuesta correcta: varios combinan rasgos que apuntan a modelos distintos, de modo que usted debe decidir qué criterio pesa más y justificarlo con datos del texto. La tabla «Datos clave» al final de cada escenario resume la información principal, pero no reemplaza la lectura del texto, donde están las tensiones que debe analizar.

\subseccion{Escenario 1. Control de temperatura y nivel en una extractora de aceite de palma}

La Extractora Río Lindo (organización ficticia creada para esta actividad) procesa fruta de palma africana en el cantón Buena Fe, provincia de Los Ríos, y recibe en temporada alta unas 180 toneladas de fruta por día. Hoy los operadores controlan a mano las válvulas de vapor de los tres esterilizadores y el nivel de los dos tanques de aceite crudo, guiándose por termómetros y reglillas de nivel. En el último año la planta registró dos desbordes de tanque y un operador sufrió quemaduras leves al abrir una compuerta con presión residual. La gerencia contrató el desarrollo de un sistema embebido que regule la temperatura y la presión de esterilización, controle el nivel de los tanques, active alarmas y bloquee la apertura de compuertas cuando exista presión.

El jefe de planta, Ing. Marcelo Andrade (persona ficticia), conoce el proceso a fondo, pero solo puede reunirse con el equipo dos mañanas por semana; los operadores trabajan en tres turnos y el responsable de seguridad industrial debe aprobar cada función de protección. Las reglas de seguridad (temperaturas y presiones máximas, tiempos de bloqueo) están definidas por el fabricante de los esterilizadores y casi no cambiarán. En cambio, el proveedor de sensores y válvulas todavía no ha sido elegido y la compra se decidirá en el tercer mes, por lo que se desconoce cómo se comunicarán esos equipos con el controlador. Los reportes de producción y la pantalla para los operadores tampoco están claros: el gerente pide algo «parecido a lo que vio en otra planta». La aseguradora de la empresa exige, antes de la puesta en marcha, un documento de requisitos de seguridad aprobado y evidencia de pruebas de cada función de protección. El controlador debe funcionar en la red local de la planta, sin conexión a internet y con memoria limitada.

\needspace{8\baselineskip}
\noindent\begin{tabularx}{\linewidth}{|L{0.30\linewidth}|Y|}
	\hline
	\thdr{Datos clave} & \thdr{Escenario 1}\\
	\hline
	Equipo & 5 personas (2 de software embebido, 1 electrónico, 1 analista de requisitos, 1 encargado de pruebas)\\
	\hline
	\rowcolor{rowgray} Plazo & 10 meses, con puesta en marcha antes de la próxima temporada alta\\
	\hline
	Disponibilidad del cliente & Baja a media: jefe de planta dos mañanas por semana; operadores en turnos rotativos\\
	\hline
	\rowcolor{rowgray} Volatilidad esperada & Baja en funciones de seguridad; alta en interfaz de operador, reportes y comunicación con sensores aún no comprados\\
	\hline
	Criticidad & Muy alta: seguridad física de operadores, desbordes y pérdidas de producción\\
	\hline
	\rowcolor{rowgray} Marco normativo o contractual & Exigencias de la aseguradora (requisitos de seguridad aprobados y pruebas documentadas) y especificaciones del fabricante de los esterilizadores\\
	\hline
\end{tabularx}

\subseccion{Escenario 2. Aplicación móvil de pedidos para un emprendimiento de comida típica}

Sabores de la Abuela Rosa (organización ficticia creada para esta actividad) es un emprendimiento familiar de Quevedo que vende bolones, tigrillo, encebollado y tortillas de yuca. Atiende un local pequeño en el centro de la ciudad y recibe unos 70 pedidos diarios por una aplicación de mensajería; la dueña, Doña Rosa Mendoza (persona ficticia), anota los pedidos en un cuaderno y con frecuencia se equivoca de dirección o pierde mensajes en las horas pico. Quiere una aplicación móvil para que los clientes vean el menú del día, hagan pedidos, elijan retirar en el local o recibir a domicilio y paguen por transferencia; más adelante, quizá, con tarjeta.

Doña Rosa responde mensajes a cualquier hora y está feliz de probar cada versión, pero cambia de idea con facilidad: una semana pide un programa de puntos para clientes frecuentes y la siguiente prefiere promociones por día. Su hijo, que administra las redes sociales, propone funciones nuevas en cada reunión. El menú cambia casi a diario según lo que se consigue en el mercado. La aplicación guardará nombres, direcciones y números de los clientes, de modo que debe cumplir la normativa ecuatoriana de protección de datos personales, y los pagos por transferencia deben conciliarse para evitar pedidos sin pagar. El presupuesto es fijo y pequeño, y Doña Rosa quiere «tener todo listo» para las fiestas de la ciudad, dentro de 14 semanas, porque espera duplicar las ventas en esa fecha. El equipo está formado por dos desarrolladores que trabajan medio tiempo. Muchos clientes usan teléfonos Android de gama baja y datos móviles intermitentes, por lo que la aplicación debe ser liviana y tolerar cortes de conexión.

\needspace{8\baselineskip}
\noindent\begin{tabularx}{\linewidth}{|L{0.30\linewidth}|Y|}
	\hline
	\thdr{Datos clave} & \thdr{Escenario 2}\\
	\hline
	Equipo & 2 desarrolladores a medio tiempo\\
	\hline
	\rowcolor{rowgray} Plazo & 14 semanas, con fecha fija por las fiestas de la ciudad y presupuesto cerrado\\
	\hline
	Disponibilidad del cliente & Muy alta: la dueña responde a diario y prueba cada versión\\
	\hline
	\rowcolor{rowgray} Volatilidad esperada & Muy alta: menú diario, ideas nuevas de la dueña y de su hijo en cada reunión\\
	\hline
	Criticidad & Media: pagos por transferencia y datos personales de clientes; sin riesgo para la vida\\
	\hline
	\rowcolor{rowgray} Marco normativo o contractual & Normativa ecuatoriana de protección de datos personales; acuerdo informal de precio fijo\\
	\hline
\end{tabularx}

\subseccion{Escenario 3. Sistema de trámites municipales contratado por contratación pública}

El GAD Municipal del cantón San Isidro del Baba (organización ficticia creada para esta actividad), en la provincia de Los Ríos, atiende cada año cerca de 9\,000 trámites de permisos de funcionamiento, patentes y certificados de no adeudar. Hoy los ciudadanos deben acudir hasta tres veces a ventanilla y los técnicos revisan carpetas en papel. El municipio adjudicó, mediante un proceso de contratación pública, el desarrollo de un sistema web de trámites en línea. Los pliegos fijan tres entregables con fecha y pago asociado: la especificación de requisitos aprobada (mes 2), el sistema instalado y probado (mes 6) y la capacitación con el acta de entrega-recepción definitiva (mes 7). Cada día de retraso genera multa, y cualquier cambio de alcance que aumente el costo requiere un trámite formal adicional que puede tardar semanas.

El administrador del contrato, Lcdo. Fabián Zambrano (persona ficticia), se reúne con el equipo una vez por semana y es el único canal oficial; los técnicos de ventanilla y la jefa de rentas, que serán los usuarios principales, solo participan cuando él los convoca. En las primeras reuniones la jefa de rentas ya cambió dos veces el flujo de aprobación de patentes, y el concejo municipal discute una nueva ordenanza de tasas que podría aprobarse en el cuarto mes, lo que modificaría los cálculos de cobro. El sistema manejará datos personales de los ciudadanos y el cobro de tasas, por lo que los errores afectan dinero público y pueden generar reclamos. El equipo de la empresa contratista es de 6 personas. El sistema debe instalarse en los servidores propios del municipio y leer la información catastral de una base de datos antigua que nadie ha documentado.

\needspace{8\baselineskip}
\noindent\begin{tabularx}{\linewidth}{|L{0.30\linewidth}|Y|}
	\hline
	\thdr{Datos clave} & \thdr{Escenario 3}\\
	\hline
	Equipo & 6 personas de la empresa contratista\\
	\hline
	\rowcolor{rowgray} Plazo & 7 meses, con tres entregables de fecha fija y multas por retraso\\
	\hline
	Disponibilidad del cliente & Media a baja: administrador del contrato una vez por semana; usuarios finales solo cuando él los convoca\\
	\hline
	\rowcolor{rowgray} Volatilidad esperada & Media a alta: cambios en el flujo de aprobación y posible nueva ordenanza de tasas en el mes 4\\
	\hline
	Criticidad & Alta: dinero público, datos personales de ciudadanos y posibles reclamos\\
	\hline
	\rowcolor{rowgray} Marco normativo o contractual & Contratación pública con pliegos, entregables y pagos fijos; cambios de alcance solo mediante trámite formal; ordenanzas municipales\\
	\hline
\end{tabularx}

\subseccion{Escenario 4. Plataforma web para una cooperativa de productores de cacao con socios dispersos}

La Cooperativa Agroproductiva Frutos del Estero (organización ficticia creada para esta actividad) agrupa a 640 pequeños productores de cacao de recintos rurales de Mocache, Ventanas y Valencia. La cooperativa acopia el cacao en un centro ubicado en Mocache, lo vende a exportadores y paga a cada socio según el peso, la humedad y la calidad de su entrega. Hoy las entregas se anotan en tarjetas de papel, las liquidaciones se calculan en una hoja de cálculo y los socios se quejan de que no entienden cómo se calculó su pago ni cuánto deben de los pequeños créditos para insumos. La directiva quiere una plataforma web donde los socios consulten entregas, liquidaciones y saldos de crédito, y donde la cooperativa publique avisos y convocatorias a asambleas.

La gerente, Ing. Patricia Vera (persona ficticia), está disponible casi a diario, pero no puede decidir sola: los cambios importantes los aprueba la directiva en su reunión mensual. Los socios viven lejos y solo se les puede consultar los días de acopio o en las asambleas trimestrales; muchos tienen poca experiencia con tecnología. Las reglas de liquidación están en el estatuto y son estables, pero los compradores exportadores cambian cada temporada los requisitos de trazabilidad y certificación del cacao, y la plataforma deberá registrarlos. Se manejará dinero de los socios y datos personales, y la cooperativa está sujeta al control del organismo que supervisa a la economía popular y solidaria. El proyecto se financia con una donación que exige informes semestrales de avance con funciones terminadas. El equipo es de 4 personas y el plazo es de 12 meses. En el centro de acopio la conexión a internet es inestable, por lo que el registro de entregas debe funcionar sin conexión y sincronizarse después.

\needspace{8\baselineskip}
\noindent\begin{tabularx}{\linewidth}{|L{0.30\linewidth}|Y|}
	\hline
	\thdr{Datos clave} & \thdr{Escenario 4}\\
	\hline
	Equipo & 4 personas\\
	\hline
	\rowcolor{rowgray} Plazo & 12 meses, con informes semestrales al donante\\
	\hline
	Disponibilidad del cliente & Mixta: gerente casi a diario; directiva una vez al mes; socios solo en días de acopio y asambleas trimestrales\\
	\hline
	\rowcolor{rowgray} Volatilidad esperada & Baja en reglas de liquidación; alta en requisitos de trazabilidad y certificación de los compradores\\
	\hline
	Criticidad & Alta: dinero y créditos de los socios, y datos personales\\
	\hline
	\rowcolor{rowgray} Marco normativo o contractual & Estatuto de la cooperativa, control del organismo supervisor de la economía popular y solidaria, condiciones del donante y normativa de protección de datos personales\\
	\hline
\end{tabularx}

\bibliographystyle{unsrtnat}
\bibliography{referencias}

\end{document}
